\begin{abstract}
Extremist communities do not simply get deplatformed and disappear; they cycle across the platform ecosystem---concentrating, raiding mainstream spaces, retreating, and re-entering. What structural features enable this pattern? I adapt the Tiebout model of citizen sorting to conceptualize platforms as quasi-governments offering public goods and communities as mobile citizen-voters choosing among governance regimes. Using an agent-based model with three simulation experiments in a $3 \times 3 \times 3$ factorial design crossing platform count, extremist proportion, and parasitism intensity, I find that the sorting mechanism is remarkably resilient to parasitism on average but breaks down under identifiable structural conditions. Three results emerge. First, platform diversification protects mainstream communities, and the premium grows with threat severity as the mechanism shifts from preference matching to target dilution. Second, coalition-based governance provides the strongest structural firewall: mediocre without extremists, it generates endogenous enclaves that sever the parasitism channel under threat. Third, the concentration--raid--retreat cycle is an emergent property of decentralized sorting, not coordinated strategy. The policy implication is that platform-system architecture---the mix of governance types and jurisdictional diversity---rather than individual moderation decisions, determines whether mainstream communities are structurally protected or exposed.

\textbf{Keywords:} social media platforms, community sorting, Tiebout model, content moderation, online extremism, agent-based modeling
\end{abstract}
